\chapter{Referencial Teórico}
\label{cap:referencial}

Este capítulo apresenta o referencial conceitual sobre o qual o
infográfico se apoia. As subseções a seguir espelham a estrutura
modular do artefato e estão ancoradas tanto nos livros-texto da
disciplina quanto em artigos seminais e material didático do
Prof.\ Koide~\cite{koide2026}.

\section{Desktop \textit{versus} Servidor}
\label{sec:ref-desktop-servidor}

Embora compartilhem o modelo de von Neumann~\cite{stallings2017},
desktop e servidor divergem em praticamente todos os subsistemas
físicos. Servidores exigem altíssima disponibilidade --- tipicamente
\textsc{sla} de 99,99\,\% (4,38 minutos de \textit{downtime} mensais)
ou 99,999\,\% (\textit{five-nines}, 26 segundos mensais) --- e
operação ininterrupta 24/7. Para isso adotam:

\begin{itemize}
    \item \textbf{Fontes redundantes \textit{hot-plug}} (configuração
          N+1 ou 2N), onde a queima de uma unidade não interrompe a
          operação;
    \item \textbf{Arquitetura multi-socket} (2, 4 ou 8 processadores
          numa mesma placa), com interconexão coerente ponto-a-ponto
          --- \textsc{qpi}/\textsc{upi} na Intel~\cite{intel2009qpi},
          HyperTransport/Infinity Fabric na AMD;
    \item \textbf{Memória \textsc{ecc} registrada} (\textsc{rdimm} ou
          \textsc{lrdimm}), capaz de detectar e corrigir erros de bit
          em tempo real;
    \item \textbf{Form factor padronizado em rack} EIA-310 (unidade
          básica \textsc{1u} $=$ 44,45\,mm), permitindo densidade e
          padronização em \textit{datacenters};
    \item \textbf{Gerenciamento \textit{out-of-band}} por
          microcontroladores dedicados como iDRAC (Dell), iLO (HPE)
          ou \textsc{ipmi} genérico, com endereço IP próprio e
          independente do sistema operacional principal;
    \item \textbf{Redes de alto desempenho} (10\,\textsc{gbe},
          25\,\textsc{gbe}, 100\,\textsc{gbe} ou \textsc{InfiniBand})
          com \textit{bonding} \textsc{lacp} para tolerância a
          falhas e agregação de banda.
\end{itemize}

O surgimento de cargas de trabalho de inteligência artificial
generativa redesenhou o conceito de ``servidor'' nos últimos cinco
anos. A planta passou a ser dominada por aceleradores
--- \textsc{gpu}s, \textsc{tpu}s, \textsc{npu}s --- com interconexões
dedicadas (\textsc{nv}Link, \textsc{cxl}) que ultrapassam o
\textsc{pci}e em ordens de grandeza~\cite{nvidia2023grace}.

\section{RAID --- Redundant Array of Independent Disks}
\label{sec:ref-raid}

A proposta original de Patterson, Gibson e Katz~\cite{patterson1988}
em 1988 introduziu a ideia de agrupar múltiplos discos baratos para
obter, ao mesmo tempo, desempenho e tolerância a falhas
\textit{barata}. Os principais níveis adotados na indústria são:

\begin{description}
    \item[RAID 0 (Striping)] divide blocos de dados sequencialmente
          entre $N$ discos; capacidade útil $N \cdot D$, leitura e
          escrita paralelas, \emph{zero} tolerância a falhas.
    \item[RAID 1 (Mirroring)] mantém cópia idêntica em dois discos;
          capacidade útil $D$, sobrevive à perda de $N-1$ unidades
          de um grupo espelhado.
    \item[RAID 5 (Paridade distribuída)] espalha um bloco de paridade
          $P = A \oplus B \oplus C$ entre os discos; capacidade útil
          $(N-1) \cdot D$, sobrevive à perda de \emph{um} disco,
          impõe penalidade de escrita 4$\times$ (\textit{read-modify-write}).
    \item[RAID 6] usa dois blocos de paridade independentes (P e Q);
          tolera \emph{dois} discos quebrados, com penalidade 6$\times$.
    \item[RAID 10] combina espelhamento + \textit{striping}; alta
          performance e tolerância, custa metade da capacidade bruta.
\end{description}

\subsection{Aplicação da Lei de Amdahl}

O argumento de Amdahl~\cite{amdahl1967} delimita o ganho máximo de
qualquer otimização paralela: o ganho total é limitado pela parcela
serial do trabalho. A formulação moderna em equação, derivada
posteriormente do argumento qualitativo original de 1967, expressa-se
como --- sendo $P$ a fração paralelizável do trabalho e $N$ o número
de unidades de execução:

\begin{equation}
\mathrm{Speedup}(N) = \frac{1}{(1-P) + \dfrac{P}{N}}
\label{eq:amdahl}
\end{equation}

Quando $N \to \infty$, $\mathrm{Speedup} \to 1/(1-P)$. No exemplo do
servidor \textsc{www} discutido pelo Prof.\ Koide na Aula~9
(disco 10\,000\,\textsc{rpm}, \textit{seek} médio de 8\,\textsc{ms},
páginas de 8\,\textsc{kb}), aproximadamente 99,3\,\% do tempo de
resposta é latência rotacional --- portanto serial --- e apenas 0,7\,\%
é transferência paralelizável. Substituindo na Equação~\ref{eq:amdahl}
com $P = 0{,}007$, o \textit{speedup} máximo teórico é
$1/(1-0{,}007) \approx 1{,}007\times$. Conclusão pedagógica: paralelizar
discos para esse perfil não ajuda; o vetor de otimização é
\emph{trocar a parte serial} --- substituir o \textsc{hdd} por
\textsc{ssd}.

\section{Memória ECC e o código de Hamming}
\label{sec:ref-ecc}

Em servidor 24/7, \textit{soft errors} de memória são
estatisticamente certos. Partículas subatômicas (nêutrons cósmicos
secundários e partículas $\alpha$ provenientes de contaminação de
solda) podem inverter um bit de \textsc{dram} sem aviso. O estudo
de campo de Schroeder, Pinheiro e Weber~\cite{schroeder2009},
realizado durante 2,5 anos na frota do Google, registrou taxas de
\textbf{25\,000 a 70\,000 FIT por Mbit} (erros por bilhão de
horas-dispositivo por Mbit) e mostrou que \textbf{mais de 8\,\% dos
\textsc{dimm}s} exibem ao menos um erro corrigível por ano ---
contrariando a hipótese anterior de erros raros. O estudo concluiu
ainda que os erros são predominantemente \textit{hard errors}
(defeitos de hardware), e não \textit{soft errors} transitórios.

A correção em tempo real é feita por um código corretor proposto
originalmente por Hamming~\cite{hamming1950} em 1950. Para 8 bits de
dado, o código clássico Hamming(12,8) anexa quatro bits de paridade
nas posições $2^k$ ($k=0,1,2,3$) e oferece \emph{single error
correction} (SEC). A variante \textsc{secded} (\emph{Single Error
Correction, Double Error Detection}) --- padrão em \textsc{dimm}
\textsc{ecc} registrado --- estende o código para (13,8) adicionando
um \emph{bit de paridade global}, capaz de detectar (mas não corrigir)
ocorrências de dois erros simultâneos. Cada bit de paridade cobre um
subconjunto de posições determinado por seu bit correspondente em
binário; na leitura, o controlador recalcula as paridades e obtém um
\emph{síndrome} cuja representação decimal é exatamente a posição do
bit corrompido. O hardware inverte o bit indicado e entrega o dado
correto à CPU sem qualquer participação do sistema operacional.

A norma JEDEC \textsc{jesd79-5}~\cite{jedec2025ddr5} tornou
obrigatória a presença de \textsc{ecc} \textit{on-die} na
\textsc{ddr5}, mas é importante notar: trata-se de proteção
\emph{interna} à célula, não cobrindo o \textit{datapath}
\textsc{dimm}~$\to$~CPU. Somente \textsc{rdimm}/\textsc{lrdimm} com
\textsc{ecc} de fato protegem todo o canal.

\section{Hierarquia de memória, cache e NUMA}
\label{sec:ref-hierarquia}

Patterson e Hennessy~\cite{patterson2017} formalizam o princípio:
nenhuma tecnologia de memória é simultaneamente rápida, grande e
barata. A solução é dispor camadas em pirâmide:
registradores (sub-nanossegundo) $\to$ cache L1/L2/L3
(nanossegundos) $\to$ \textsc{ram} (dezenas de nanossegundos) $\to$
\textsc{ssd} (microssegundos) $\to$ \textsc{hdd} (milissegundos).
A eficiência da pirâmide depende de duas formas de localidade:
\textbf{temporal} (dados acessados recentemente serão acessados
de novo) e \textbf{espacial} (vizinhos de um dado acessado serão
acessados em breve).

\subsection{Organização da cache}

Caches modernos são \textit{set-associative}: o endereço seleciona um
\textit{set} e a linha pode ocupar qualquer um dos $k$ \textit{ways}.
Um Intel Core i7-5960X, processador de referência citado no
material de aula~\cite{koide2026}, organiza-se assim:

\begin{center}
\begin{tabular}{lrrrr}
\toprule
Nível & Tamanho & Associatividade & Linha & Latência \\
\midrule
L1 (dados/instruções) & 32 KB     & 8-way  & 64 B & 4 ciclos \\
L2                    & 256 KB    & 8-way  & 64 B & 12 ciclos \\
L3 (\emph{LLC})       & 20 MB     & 20-way & 64 B & 38 ciclos \\
\bottomrule
\end{tabular}
\end{center}

A \textbf{linha de cache de 64 bytes} é a unidade indivisível: ler
um byte traz consigo 63 vizinhos. L1 é privada de cada \emph{core},
L3 é compartilhada por todos os \emph{cores} do mesmo \emph{socket}.

\subsection{NUMA --- Non-Uniform Memory Access}

Em servidor multi-socket, cada CPU possui seus pentes próprios de
RAM. O acesso ao banco local é rápido ($\approx 80$\,\textsc{ns}),
mas o acesso à RAM do socket vizinho passa pelo link
\textsc{upi}/Infinity Fabric e custa cerca de 140\,\textsc{ns}~\cite{stallings2017}.
Essa não-uniformidade obriga o escalonador a manter cada
\emph{thread} próxima da memória que ela toca --- conceito de
\textit{NUMA affinity} suportado por \texttt{numactl} (Linux) e
funções equivalentes da \textsc{api} Windows.

\subsection{Coerência: MESI, MESIF, MOESI}

Quando múltiplos \emph{cores} têm cópias da mesma linha de cache, é
preciso garantir consistência. O protocolo \textsc{mesi} define
quatro estados por linha: \emph{Modified, Exclusive, Shared, Invalid}.
A Intel acrescenta \emph{Forward} (\textsc{mesif}) e a AMD acrescenta
\emph{Owned} (\textsc{moesi}). Toda escrita em linha
\emph{Shared} dispara invalidações nos outros \emph{cores}; essa
operação cruza o link de coerência e custa centenas de
ciclos --- o conhecido \textit{false sharing}.

\subsection{Paginação, TLB e \textit{huge pages}}

A \textsc{mmu} traduz endereços virtuais para físicos via tabelas de
página percorridas em 4 níveis na arquitetura x86-64. A \textsc{tlb}
(\emph{Translation Lookaside Buffer}) cacheia as traduções
recentes; quando há \textit{miss}, o \textit{page walk} custa até
quatro acessos à \textsc{ram}~\cite{silberschatz2015}. Para reduzir
pressão na \textsc{tlb}, sistemas modernos suportam páginas de
2\,\textsc{mb} (\textit{huge pages}) e 1\,\textsc{gb}
(\textit{gigantic pages}), cada entrada cobrindo 512$\times$ ou
262\,144$\times$ mais memória.

\section{Virtualização e hipervisores}
\label{sec:ref-virtualizacao}

Popek e Goldberg~\cite{popek1974} estabeleceram, em 1974, três
condições para que uma arquitetura seja virtualizável:
\textbf{equivalência} (a VM deve se comportar como hardware nativo),
\textbf{eficiência} (a maior parte das instruções deve executar sem
intervenção do hipervisor) e \textbf{controle de recursos}
(o hipervisor mantém posse exclusiva dos recursos físicos). Esses
critérios não eram naturalmente satisfeitos pela arquitetura
x86 original; somente em 2005 (Intel \textsc{vt-x}) e 2006 (AMD-V) as
extensões de hardware introduziram o estado \textsc{vmx} (informalmente
chamado de Ring~$-1$) que coloca o hipervisor abaixo do
Ring~0~\cite{russinovich2017}.

\subsection{Tipo 1 \textit{vs.}\ Tipo 2}

\begin{itemize}
    \item \textbf{Tipo 1 (\textit{bare metal})}: o hipervisor é o
          próprio sistema operacional raiz, instalado diretamente
          sobre o hardware. Exemplos: VMware ESXi, KVM (kernel
          Linux + \textsc{qemu}), Hyper-V Server, Proxmox VE, Citrix
          XenServer. Performance próxima da nativa.
    \item \textbf{Tipo 2 (\textit{hosted})}: o hipervisor é um
          programa rodando sobre um sistema operacional convencional.
          Exemplos: VirtualBox (citado pelo Prof.\ Koide na Aula~12),
          VMware Workstation, Parallels Desktop. Performance reduzida
          pelo \textit{overhead} do \textsc{so} hospedeiro.
\end{itemize}

O Windows Subsystem for Linux 2 (\textsc{wsl}~2) ilustra o caso
limítrofe: é um \emph{kernel} Linux verdadeiro rodando numa VM
\textit{utility} leve gerenciada pelo Hyper-V, com \textit{boot}
quase instantâneo e integração transparente com o filesystem do
Windows via 9P.

\subsection{Contêineres: virtualização ao nível do SO}

Diferentemente da VM, o contêiner compartilha o \emph{kernel} do
\textit{host} --- isolamento é obtido por \textit{namespaces}~\cite{man7namespaces}
(PID, mount, network, IPC, UTS, user) e \textit{cgroups}~\cite{kernelcgroups}
do Linux. O resultado: \textit{boot} em $\approx
100$\,\textsc{ms}, \textit{footprint} de \textsc{mb}s, densidade
acima de 1\,000 contêineres por \emph{host}.
Orquestradores como Kubernetes escalam essa unidade para clusters de
milhares de nós.

\subsection{Escalonadores}

O \emph{Completely Fair Scheduler} (\textsc{cfs}) do Linux, introduzido
no \emph{kernel} 2.6.23 (2007), mantém todas as \texttt{task\_struct}
numa árvore rubro-negra ordenada por \texttt{vruntime} --- tempo
virtual de execução ponderado pela prioridade \emph{nice}. A cada
decisão, escolhe-se a tarefa mais à esquerda; complexidade
$\mathcal{O}(\log n)$~\cite{love2010}. O
Windows usa escalonador \emph{multilevel feedback queue} baseado em
quantum e prioridades 0--31, com a maior parte da política
implementada no \texttt{ntoskrnl.exe}~\cite{russinovich2017}.

\section{Síntese: \textit{stack} completo do servidor}
\label{sec:ref-stack}

Os cinco subsistemas anteriores não são compartimentos isolados.
Numa requisição HTTP simples a um \textit{e-commerce}, todos eles
participam: o \textit{load balancer} entrega o pacote pela rede
(camada~\ref{sec:ref-desktop-servidor}); o \textit{ingress
controller} no Kubernetes despacha para um \emph{pod}
(virtualização, \ref{sec:ref-virtualizacao}); o processo executa em
v\textsc{cpu}s mapeados em \emph{cores} físicos com afinidade
\textsc{numa} (\ref{sec:ref-hierarquia}); buffers são alocados em
\textsc{ram} \textsc{ecc} (\ref{sec:ref-ecc}); páginas são lidas de
um \textsc{ssd} integrado a \textsc{raid}~10 (\ref{sec:ref-raid}); e
a resposta retorna percorrendo a mesma cadeia em sentido inverso.

A última seção do infográfico, descrita no
Capítulo~\ref{cap:desenvolvimento}, materializa essa visão
integradora.
